\section{Discussion and Limitations}
The prototype links physical sensing and guided checks through evidence-handling
contracts. State-transition and API tests verify symptom retention, correction
history, and report isolation under constructed interactions. They establish
implementation behavior; participant records separately describe target response,
healthy measurements, and deployment.

These observations expose the remaining application question: can guided checks
reject ordinary-activity triggers while retaining abnormal signs? The recorded
0.480 false triggers per reference-negative online hour represents a practical
burden. Healthy-check positivity cannot establish a low final-alert rate because
protocols, baselines, and output definitions differ. A frozen-configuration study
linking every trigger, guided outcome, final alert, and missed target would test
the workflow's benefit directly.

Historical configuration contrasts remain descriptive because comparator
settings are unavailable. The complete-first rule comparison isolates a policy
difference, while API traces verify synthetic interactions through reporting and
storage. Neither reconstructs the historical symptom and history failures;
their causes require original input sequences.

Component evidence also has specific limits: MDSC sensitivity varies across six
test speakers; Face reference uncertainty is undefined; Eyes exported values lack
a unique runtime mapping. Clinical sensitivity requires prospective clinical
references, and predictive fusion requires same-session multimodal data with
participant-level separation \cite{collins2024tripodai}. The two-person pilot
supports physical deployment but records outages, data loss, and a restart.
Continuous 24-h operation and external alert delivery remain untested.

\section{Conclusion}
PiBE-FAST integrates passive sensing, guided BE-FAST checks, and evidence-preserving
interactions for stroke warning-sign assessment. Software verification and
participant observations establish a working prototype and expose a remaining
false-trigger burden. Reliable, low-burden warning requires linked final-alert
outcomes and independent clinical validation.
